\addtocontents{toc}{\protect\Needspace{3\baselineskip}}
\originalchapter{2007 年第一次正式测验}
五题取自 2007 年 3 月 21 日 Test 1, 考试时长 90 分钟, 允许个人两页笔记. 解答为本稿独立推导. 锥形空间题已按原卷第 3 页附图核对并用 TikZ 重绘; 原题未规定点的精确半径, 本稿也不补造数值条件.

\src{8.251 Test 1, Q1, 10 分.}
\begin{exercise}
四维时空 $(ct,x^1,x^2,x^3)$ 中两静止 D2 膜分别位于 $x^3=0$ 和 $x^2=0$. 画示意. 开弦两端 $\sigma=0,\sigma_1$ 分别在第一、第二膜上, 列出每端四坐标的共八条自由或 Dirichlet 边界条件及固定值.
\end{exercise}
\begin{solution}
以下取共形规范, 自由方向的边界动量流 $P^\sigma_\mu=0$ 可写成 $X'^\mu=0$; 未固定规范时应保留动量流条件. 第一端 $X'^0=X'^1=X'^2=0,X^3=0$; 第二端 $X'^0=X'^1=X'^3=0,X^2=0$. 时间与共有 $x^1$ 方向都自由, 每个膜各有一个法向固定. 示意画垂直公共 $x^1$ 的截面, 第一膜为 $x^2$ 轴、第二膜为 $x^3$ 轴:
\begin{center}\begin{tikzpicture}
\draw[-{Stealth},thick](-2,0)--(2,0)node[right]{$x^2$};\draw[-{Stealth},thick](0,-1.3)--(0,2)node[above]{$x^3$};
\draw[structurecolor,thick](1.4,0)..controls(1,1)and(.4,1.5)..(0,1.6);
\fill(1.4,0)circle(2pt)node[below]{$\sigma=0$};\fill(0,1.6)circle(2pt)node[left]{$\sigma_1$};
\end{tikzpicture}\end{center}
这只是端点条件示意, 不把任意所画曲线当作满足全部弦约束的具体动力学解.
\end{solution}

\src{8.251 Test 1 Q3, (a)--(b), 10 分.}
\begin{exercise}
取历史题面真空质量密度 $\rho_v=7.6\times10^{-27}\,\mathrm{kg\,m^{-3}}$. 仅用 $\rho_v,\hbar,c$ 构造长度 $\ell_v=\rho_v^\alpha\hbar^\beta c^\gamma$.
(a) 求三指数; (b) 取 $\hbar=1.055\times10^{-34}\,\mathrm{J\,s}$、$c=2.998\times10^8\,\mathrm{m\,s^{-1}}$, 给微米数值.
\end{exercise}
\begin{solution}
质量、时间、长度次数方程分别为 $\alpha+\beta=0$, $-\beta-\gamma=0$, $-3\alpha+2\beta+\gamma=1$. 解得 $\alpha=-1/4,\beta=1/4,\gamma=-1/4$, 故 $\ell_v=(\hbar/(\rho_vc))^{1/4}\simeq82.49\,\mu\mathrm m$. 这只是一项量纲尺度; 题面没有给能使该尺度必然对应某个新引力效应的动力学模型.
\end{solution}

\src{8.251 Test 1 Q4, (a)--(b), 20 分.}
\begin{exercise}
曲面 $z=h(x,y)$ 嵌入三维 Euclid 空间, 诱导度规 $g_{ij}=\partial_i\boldsymbol X\cdot\partial_j\boldsymbol X$.
(a) 用 $h$ 的导数写四分量和 $g_{ij}=\delta_{ij}+\cdots$.
(b) 完全化简面积积分.
\end{exercise}
\begin{solution}
切向量 $X_x=(1,0,h_x),X_y=(0,1,h_y)$, 故 $g_{ij}=\delta_{ij}+h_ih_j$, 分量为 $1+h_x^2,h_xh_y,h_xh_y,1+h_y^2$. 行列式 $(1+h_x^2)(1+h_y^2)-h_x^2h_y^2=1+h_x^2+h_y^2$, 因而 $A=\int\dd x\dd y\sqrt{1+|\nabla h|^2}$. 该面积式针对单值图面参数域, 不覆盖无法写成全局高度函数的任意闭曲面.
\end{solution}

\src{8.251 Test 1 Q5, (a)--(d), 40 分.}
\begin{exercise}
两平行 D2 膜位于 $z=0,a>0$. 弦参数 $0\le\sigma\le L$, $\boldsymbol x=(x,y)=\frac12[F(ct+\sigma)+G(ct-\sigma)]$, $z=\frac12[f(ct+\sigma)+g(ct-\sigma)]$.
(a) 用两端条件求 $G,g$ 与 $F,f$ 的关系及周期条件, 写简化解.
(b) 施加 $|\boldsymbol X'\pm\dot{\boldsymbol X}/c|^2=1$ 求约束.
若 $\sigma=0$ 端点以固定角速度 $\omega>0$ 作半径 $R$ 的圆, 且选择与 $\omega,a$ 相容的最低能解:
(c) 求 $F,L$; (d) 求 $f$ 与 $\omega R/c$.
\end{exercise}
\begin{solution}
平行方向 Neumann 条件给 $G'=F'$, 平移积分常数后可取 $G=F$, 且 $F'(u+2L)=F'(u)$. 法向 $z(t,0)=0$ 给 $g=-f$, 另一端给 $f(u+2L)-f(u)=2a$. 因而 $\boldsymbol x=\frac12[F(ct+\sigma)+F(ct-\sigma)]$, $z=\frac12[f(ct+\sigma)-f(ct-\sigma)]$. 约束分别要求 $|F'(u)|^2+f'(u)^2=1$.

端点轨迹等于 $F(ct)$, 所以 $F(u)=R(\cos(\omega u/c),\sin(\omega u/c))$. 其导数周期要求 $L=m\pi c/\omega$, $m$ 为正整数. 为使端点高度差恒定且满足约束, 取 $f'(u)=a/L$ (连续低能支), $f(u)=au/L+f_0$. 所以
\[
\frac{\omega R}c=\sqrt{1-\left(\frac{a\omega}{m\pi c}\right)^2},\qquad
E=T_0L=\frac{m\pi cT_0}{\omega}.
\]
若 $a\omega<\pi c$, 最低能 $m=1$, 得原考题通常期待的基本振动解. 若不满足这一额外范围, 不能产生虚半径; 须取最小正整数 $m>a\omega/(\pi c)$, 才有非零圆轨迹. 等号只有 $R=0$ 的退化情形, 不再是正半径圆. 题面没有排除高分支, 本解把参数条件明示出来.
\end{solution}
